% from ega0-16-fr.tex
\oldpage[0\textsubscript{IV-1}]{35}
il suffit de prendre pour $A$ un anneau local intègre de dimension $\geq1$~; on a
$\operatorname{prof}(A)\geq1$ mais $\operatorname{prof}(K)=0$ pour le corps
des fractions $K$ de $A$. Pour un exemple du second cas, considérons un
anneau local noethérien intègre $A_0$ de dimension $\geq2$, d'idéal maximal
$\mathfrak{m}_0$ et soit $k=A_0/\mathfrak{m}_0$ son corps résiduel~; soit
$A=A_0\oplus\mathfrak{J}$ l'extension triviale de $A_0$ par le $A_0$-module
$k$ (18.2.3), l'idéal $\mathfrak{J}$ étant $A_0$-isomorphe à $k$ (de sorte
que la multiplication dans $A$ est donnée par
$(x,z)(x',z')=(xx',xz'+x'z)$). Il est clair que $\mathfrak{J}$ est le
nilradical de $A$ et $A/\mathfrak{J}$ est isomorphe à $A_0$, de sorte que
tout idéal premier $\mathfrak{p}$ de $A$ est de la forme
$\mathfrak{p}_0\oplus\mathfrak{J}$, où $\mathfrak{p}_0$ est un idéal premier
de $A_0$~; en particulier $\mathfrak{m}=\mathfrak{m}_0\oplus\mathfrak{J}$
est l'unique idéal maximal de $A$. Comme tout élément de $\mathfrak{m}_0$
annule $\mathfrak{J}$, on a
$\operatorname{prof}(A)=\operatorname{prof}(A_{\mathfrak{m}})=0$~; au
contraire, si $\mathfrak{p}_0\neq\mathfrak{m}_0$, les éléments de
$\mathfrak{J}$ sont annulés par ceux de
$\mathfrak{m}_0-\mathfrak{p}_0$, donc $A_{\mathfrak{p}}$ est canoniquement
isomorphe à $(A_0)_{\mathfrak{p}_0}$, et comme $A_0$ est intègre,
$\operatorname{prof}(A_{\mathfrak{p}})
=\operatorname{prof}((A_0)_{\mathfrak{p}_0})\geq1$ si
$\mathfrak{p}_0\neq0$.
\end{env}

\begin{corollary}[16.4.7]
\label{0.16.4.7-fr}
Pour qu'un anneau local noethérien réduit $A$ soit de profondeur $0$, il faut
et il suffit que $A$ soit un corps.
\end{corollary}

\begin{proof}
En effet, comme l'intersection des idéaux premiers minimaux de $A$ est $0$,
$A$ n'a pas d'idéaux premiers associés immergés~; en vertu de (16.4.6, (i)),
si $\operatorname{prof}(A)=0$, l'idéal maximal $\mathfrak{m}$ de $A$ est
aussi idéal premier minimal, donc c'est le seul idéal premier de $A$, et
comme $A$ est réduit, $\mathfrak{m}=0$.
\end{proof}

\begin{proposition}[16.4.8]
\label{0.16.4.8-fr}
Soient $A$, $B$ deux anneaux locaux noethériens,
$\rho:A\to B$ un homomorphisme local, $M$ un $B$-module tel que
$M_{[\rho]}$ soit un $A$-module de type fini. Alors on a
\begin{equation}
\label{0.16.4.8.1-fr}
\tag{16.4.8.1}
  \operatorname{prof}_A(M_{[\rho]})=\operatorname{prof}_B(M).
\end{equation}
\end{proposition}

\begin{proof}
On peut se borner au cas $M\neq0$~; supposons que
$\operatorname{prof}_A(M_{[\rho]})=n$. Si $(x_i)_{1\leq i\leq n}$ est une
suite $(M_{[\rho]})$-régulière maximale formée d'éléments de l'idéal maximal
de $A$, les $\rho(x_i)$ constituent une suite $M$-régulière formée d'éléments
de l'idéal maximal de $B$, et on a donc, en posant
$N=M/(\rho(x_1)M+\cdots+\rho(x_n)M)$,
$\operatorname{prof}_B(N)=\operatorname{prof}_B(M)-n$ (16.4.6, (ii))~; de
même
$N_{[\rho]}=M_{[\rho]}/(x_1M_{[\rho]}+\cdots+x_nM_{[\rho]})$ et
$\operatorname{prof}_A(N_{[\rho]})=0$~; on est ainsi ramené à démontrer la
proposition lorsque $n=0$. Soit $k$ le corps résiduel de $A$~; comme $k$ est
un quotient de $A$, le $A$-module
$P=\operatorname{Hom}_A(k,M_{[\rho]})$ est un sous-module de
$\operatorname{Hom}_A(A,M_{[\rho]})=M_{[\rho]}$ et l'hypothèse entraîne que
$P\neq0$ (16.4.6, (i))~; en outre, $P$ est un $A$-module de type fini
(puisque $M_{[\rho]}$ est un $A$-module de type fini), donc un $k$-espace
vectoriel de type fini, étant annulé par l'idéal maximal de $A$, et finalement
un $A$-module de longueur finie. D'autre part, $P$ est aussi un sous-module du
$B$-module $M$, et puisqu'il est de longueur finie en tant que $A$-module, il
l'est \emph{a fortiori} en tant que $B$-module. Comme il est $\neq0$, il
contient un sous-$B$-module simple, c'est-à-dire isomorphe au corps résiduel de
$B$~; on conclut alors de (16.4.6, (i)) que l'on a bien
$\operatorname{prof}_B(M)=0$.
\end{proof}

\begin{definition}[16.4.9]
\label{0.16.4.9-fr}
Soient $A$ un anneau local noethérien, $M$ un $A$-module de type fini. Si
$M\neq0$, on appelle \emph{coprofondeur} de $M$ et on note
$\operatorname{coprof}_A(M)$ ou $\operatorname{coprof}(M)$, l'entier fini
$\dim_A(M)-\operatorname{prof}_A(M)\geq0$ (16.4.5.1). Si $M=0$, on pose
$\operatorname{coprof}(M)=0$.

On a $\operatorname{coprof}(M^k)=\operatorname{coprof}(M)$ pour tout $k>0$.
\end{definition}

\begin{proposition}[16.4.10]
\label{0.16.4.10-fr}
Soient $A$ un anneau local noethérien, $M$ un $A$-module de type fini.
